% Part: proof-theory
% Chapter: natural-deduction
% Section: sequents

\documentclass[../../../include/open-logic-section]{subfiles}

\begin{document}

\olfileid[pt]{pt}{nat}{seq}

\olsection{Dedução natural com sequentes: \Log{N2c} e~\Log{N2i}}

!!^a{proof}~$\delta$ em \Log{G1c} ou \Log{G1i} mostra que sua
!!{formula} final~$!A$ decorre de suas hipóteses abertas. Se $\Gamma$ é
o conjunto das !!{formula}s~$!B$ para as quais $!B^x$ ocorre como
hipótese aberta em~$\delta$, podemos escrever isso como $\delta \Proves
\Gamma \Sequent !A$. Isso sugere que também podemos formular um sistema
de dedução natural para sequentes. Nesse sistema, cada sequente contém a
informação sobre as hipóteses abertas das quais depende a !!{formula} à
direita, isto é, sobre quais hipóteses estão abertas na sub-!!{proof} que
termina em um sequente. Trata-se da dedução natural com sequentes, ou
dedução natural ``em estilo de sequentes'', e chamamos os sistemas
resultantes de \Log{N2c} e~\Log{N2i}.

Para aproximar a correspondência entre \Log{N1} e \Log{N2}, tomemos os
conjuntos~$\Gamma$ à esquerda como conjuntos de !!{formula}s
rotuladas~$x:!B$. Como as regras de \Log{N1} operam apenas sobre as
!!{formula}s nas premissas, isto é, as !!{formula}s finais das
sub-!!{proof}s imediatas, as regras de \Log{N2} operam apenas sobre o
sucedente do sequente. Assim, podemos chamar simplesmente de
contexto as fórmulas rotuladas à esquerda dos sequentes.

Em vez de hipóteses, as provas de \Log{N2} têm como sequentes superiores
axiomas da forma
\[x:!A \Sequent !A.\]
As regras correspondem exatamente às regras de \Log{N1} e são
apresentadas em \olref{tab:N2}.

\subfile{rules-N2}

Como precisamos admitir (tal como em \Log{N1c} e \Log{N1i}) que as
regras \Intro{\lif}, \Intro{\lnot}, \Elim{\lor} e \Elim{\lexists}
possam, mas não precisem, descarregar efetivamente hipóteses, também
admitimos aplicações corretas das regras correspondentes em \Log{N2c}
e~\Log{N2i} quando as !!{formula}s de contexto correspondentes $!A$,
$!B$ e $!A(a)$ não estão presentes no contexto da premissa pertinente.

\begin{prop}
Se $\delta$ é !!a{proof} de~$!A$ em~\Log{N1c} ou \Log{N1i}, então
existe !!a{proof}~$\delta'$ de $\Gamma \Sequent !A$ em~\Log{N2c}
(\Log{N2i}), onde $\Gamma$ é o conjunto de todos os~$x:!B$ tais que
$!B^x$ é uma hipótese aberta de~$\delta$.
\end{prop}

\begin{proof}
  Procedemos por indução em $n=\pheight{\delta}$. Se $n=0$, $\delta$ é
  uma única hipótese aberta~$!A^x$. Então $\Gamma = \{x:!A\}$ e
  $\Gamma \Sequent !A$ é um axioma de \Log{N2c} ou~\Log{N2i}.

  Se $n>0$, distinguimos casos conforme a última inferência
  em~$\delta$. Examinamos apenas um caso; os demais ficam como
  exercícios.

  Suponha que a última inferência seja $\Intro{\lif}$. Então $\delta$
  termina em
  \[
  \AxiomC{$\Discharge{!B}{x}$}
  \RightLabel{$\delta_1$}
  \DeduceC{$!C$}
  \DischargeRule{\Intro{\lif}}{x}
  \UnaryInfC{$!B \lif !C$}
  \DisplayProof
  \]
  Pela hipótese de indução, existe !!a{proof}~$\delta_1'$ em~\Log{N2c}
  (\Log{N2i}) de $\Gamma' \Sequent !C$, onde $\Gamma'$ são as
  hipóteses rotuladas abertas em~$\delta_1$. Se $!B^x$ é uma hipótese
  aberta de $\delta_1$, ela é descarregada na inferência
  \Intro{\lif}, e as hipóteses abertas de~$\delta$ são $\Gamma =
  \Gamma' \setminus \{x:!B\}$. Em outras palavras, $\delta_1'$ é
  !!a{proof} de~$x:!B, \Gamma \Sequent !C$, e podemos tomar $\delta'$
  como
  \[
  \AxiomC{}
  \RightLabel{$\delta_1'$}
  \Deduce$x:B, \Gamma \fCenter !C$
  \RightLabel{\Intro{\lif}}
  \UnaryInf$\Gamma \fCenter !B \lif !C$
  \DisplayProof
  \]
  Se $!B^x$ não é uma hipótese aberta de~$\delta_1$, então
  \Intro{\lif} não descarrega uma hipótese, e as hipóteses abertas de
  $\delta$ e $\delta_1$ são as mesmas. Como não se exige que a fórmula
  de contexto $x:!B$ esteja presente na premissa de \Intro{\lif}
  em~\Log{N2c} e~\Log{N2i}, podemos tomar $\delta'$ como
  \[
  \AxiomC{}
  \RightLabel{$\delta_1'$}
  \Deduce$\Gamma \fCenter !C$
  \RightLabel{\Intro{\lif}}
  \UnaryInf$\Gamma \fCenter !B \lif !C$
  \DisplayProof
  \]

  Os casos em que $\delta$ termina em outra regra são semelhantes.
\end{proof}

\begin{prop}
Se $\delta$ é !!a{proof} em \Log{N2c} ou \Log{N2i} de $\Gamma
\Sequent !A$, então existe !!a{proof}~$\delta'$ em \Log{N1c}
(\Log{N1i}) com !!{formula} final~$!A$ e hipóteses abertas $!B^x$ para
todo $x:!B \in \Gamma$.
\end{prop}

\begin{proof}
  Novamente, procedemos por indução na !!{height}~$n$ de~$\delta$. Se
  $n=0$, $\delta$ é um axioma, $x:!A \Sequent !A$. A
  !!{proof}~$\delta'$ é a hipótese~$!A^x$.

  Se $n>0$, distinguimos casos conforme a última inferência
  em~$\delta$. Por exemplo, se essa inferência é~$\Intro{\lif}$,
  $\delta$~termina em
  \[
  \bottomAlignProof
  \AxiomC{}
  \RightLabel{$\delta_1$}
  \Deduce$x:!B, \Gamma \fCenter !C$
  \RightLabel{\Intro{\lif}}
  \UnaryInf$\Gamma \fCenter !B \lif !C$
  \DisplayProof
  \quad\text{ ou }\quad
  \bottomAlignProof
  \AxiomC{}
  \RightLabel{$\delta_1$}
  \Deduce$\Gamma \fCenter !C$
  \RightLabel{\Intro{\lif}}
  \UnaryInf$\Gamma \fCenter !B \lif !C$
  \DisplayProof
  \]
  Pela hipótese de indução, existe !!a{proof} de~$!C$ em~\Log{N1c}
  (\Log{N1i}). No primeiro caso, $!B^x$ é uma hipótese aberta
  de~$\delta'$; no segundo, não é. Podemos tomar a !!{proof}~$\delta'$
  como
  \[
  \bottomAlignProof
  \AxiomC{$\Discharge{!B}{x}$}
  \RightLabel{$\delta_1$}
  \DeduceC{$!C$}
  \DischargeRule{\Intro{\lif}}{x}
  \UnaryInfC{$!B \lif !C$}
  \DisplayProof
  \quad\text{ ou }\quad
  \bottomAlignProof
  \AxiomC{}
  \RightLabel{$\delta_1$}
  \DeduceC{$!C$}
  \RightLabel{\Intro{\lif}}
  \UnaryInfC{$!B \lif !C$}
  \DisplayProof
  \]
  respectivamente.
\end{proof}

\end{document}
